% ===========================================================
% Chapter 16: Multi-Agent Systems and RL
% Part IV: Frontiers
% ===========================================================

\chapter{Мультиагентные системы и RL}
\label{ch:agents}

\begin{quote}
\itshape
Целое больше суммы своих частей.
\par\medskip
\upshape\raggedleft
---\,Аристотель, \textit{Метафизика}
\end{quote}

\bigskip

Все рассмотренные нами до сих пор обстановки включают одного обучающего агента, взаимодействующего
с миром, чья динамика фиксирована. Среда может быть
стохастической, частично наблюдаемой или высокоразмерной, но сама она не
\emph{обучается}. В реальном мире это допущение редко выполняется.
Финансовые рынки населены трейдерами, каждый из которых адаптируется к
стратегиям других. Автономные транспортные средства должны вести переговоры на
перекрёстках с другими транспортными средствами, одновременно пытающимися безопасно
навигировать. Команда роботов-спасателей должна разделить здание между
собой без централизованной координации. Языковые модели, развёрнутые как
помощники, всё чаще делегируют подзадачи специализированным
агентам\,---\,другим моделям, которые отвечают, проверяют или критикуют. В каждом
случае «среда» содержит \emph{других агентов}, и эти агенты
сами обучаются и адаптируются.

Мультиагентное обучение с подкреплением (MARL)\index{multi-agent reinforcement learning|textbf}
изучает, как несколько субъектов, принимающих решения, могут обучаться хорошему
поведению при совместном существовании и взаимодействии. Эта область находится на
пересечении теории игр, оптимизации и современного глубокого обучения,
и она ставит качественно иные задачи по сравнению с
одноагентным RL: среда \emph{нестационарна}, поскольку
другие агенты меняются; вознаграждения могут быть общими, конкурирующими или
чем-то средним; а пространство совместных поведений экспоненциально больше пространства
индивидуальных поведений.

В этой главе теория и практика MARL развиваются постепенно.
Раздел~\ref{sec:markov-games} вводит марковские игры\,---\,естественное
многоагентное обобщение MDP,\,---\,и определяет концепцию
равновесия Нэша. Раздел~\ref{sec:cooperative}
охватывает кооперативный MARL, включая парадигму централизованного обучения
с децентрализованным исполнением (CTDE) и алгоритмы QMIX и MAPPO.
Раздел~\ref{sec:competitive} изучает конкурентные и
смешанные обстановки, включая игры с нулевой суммой, самоигру и историю
AlphaGo и AlphaZero. Раздел~\ref{sec:communication} исследует
возникающую коммуникацию и социальные дилеммы.
Раздел~\ref{sec:marl-llms} переносит этот подход на большие языковые
модели, охватывая дебаты, многоагентную генерацию кода и agentic
workflows. Раздел~\ref{sec:marl-challenges} разбирает открытые
проблемы, а Раздел~\ref{sec:marl-summary} подводит итоги.

% ===========================================================
\section{Мультиагентные марковские игры}
\label{sec:markov-games}
% ===========================================================

\subsection{От MDP к марковским играм}

Напомним из Главы~\ref{ch:mdp}, что марковский процесс принятия решений~--- это
кортеж $(\cS, \cA, \cT, \cR, \gamma)$, в котором один агент принимает
действия, среда переходит в новое состояние, а агент
получает скалярное вознаграждение. Когда есть $N \geq 2$ агентов,
естественным расширением является \emph{марковская игра}\index{Markov game|textbf},
также называемая \emph{стохастической игрой}\index{stochastic game}.

\begin{definition}[Марковская игра]
\label{def:markov-game}
\textbf{Марковская игра} с $N$ агентами~--- это кортеж
\[
  \mathcal{G} = \bigl(\cS,\, \cA^1, \ldots, \cA^N,\, \cT,\,
                       R^1, \ldots, R^N,\, \gamma\bigr),
\]
где:
\begin{itemize}[leftmargin=*]
  \item $\cS$~--- конечное множество \textbf{состояний}, разделяемых всеми агентами.
  \item $\cA^i$~--- \textbf{множество действий} агента $i$. \textbf{Совместное пространство действий}~--- $\cA = \cA^1 \times \cdots \times \cA^N$.
  \item $\cT : \cS \times \cA \to \Delta(\cS)$~--- \textbf{ядро перехода}; следующее состояние зависит от текущего
    состояния \emph{и} совместного действия всех агентов.
  \item $R^i : \cS \times \cA \times \cS \to \R$~--- \textbf{функция вознаграждения} для агента $i$. У каждого агента своя
    функция вознаграждения, которая может зависеть от действий всех агентов.
  \item $\gamma \in [0,1)$~--- общий коэффициент дисконтирования.
\end{itemize}
\end{definition}

При $N = 1$ марковская игра сводится к стандартному MDP. Когда все
агенты разделяют одну и ту же функцию вознаграждения, $R^1 = \cdots = R^N$, игра
является \textbf{полностью кооперативной}\index{fully cooperative game}. Когда
$\sum_i R^i = 0$ при каждом переходе, она является \textbf{игрой с нулевой суммой}\index{zero-sum game}.
Общий случай, когда вознаграждения ни идентичны, ни суммируются до
нуля, называется \textbf{смешанной или игрой с общей суммой}\index{general-sum game}.

Совместная стратегия\index{joint policy} всех агентов~---
$\boldsymbol{\pi} = (\pi^1, \ldots, \pi^N)$, где каждая $\pi^i :
\cS \to \Delta(\cA^i)$~--- стохастическая стратегия агента $i$.
Цель агента $i$~--- максимизировать его собственный дисконтированный доход:
\begin{equation}
\label{eq:marl-return}
J^i(\boldsymbol{\pi}) = \E_{\boldsymbol{\pi}}\!\left[
  \sum_{t=0}^{\infty} \gamma^t R^i(s_t, \bm{a}_t, s_{t+1})
\right],
\end{equation}
где $\bm{a}_t = (a_t^1, \ldots, a_t^N)$ извлекается из произведения
стратегий $\prod_i \pi^i(\cdot \mid s_t)$.



\begin{figure}[ht]
\centering
\begin{tikzpicture}[
  state/.style={circle, draw, minimum size=1.0cm, font=\small},
  agent/.style={rectangle, draw, rounded corners, minimum width=1.8cm,
                minimum height=0.7cm, font=\small, fill=gray!12},
  action/.style={->, >=Stealth, thick},
  reward/.style={->, >=Stealth, dashed, red!70!black},
  trans/.style={->, >=Stealth, thick, blue!60!black}
]

% Состояния
\node[state] (s) at (0,0) {$s_t$};
\node[state] (sn) at (5,0) {$s_{t+1}$};

% Агенты
\node[agent] (a1) at (-2.2, 1.5) {Агент $1$};
\node[agent] (aN) at (-2.2,-1.5) {Агент $N$};
\node[font=\small] (dots) at (-2.2, 0) {$\vdots$};

% Действия в состояние
\draw[action] (a1.east) -- node[above, font=\scriptsize] {$a_t^1$}
              ([yshift=3pt]s.west);
\draw[action] (aN.east) -- node[below, font=\scriptsize] {$a_t^N$}
              ([yshift=-3pt]s.west);

% Переход
\draw[trans] (s) -- node[above, font=\small] {$\mathcal{T}(s_{t+1}\mid s_t,\bm{a}_t)$} (sn);

% Стрелки вознаграждения
\draw[reward] (sn.north) -- ++(0, 1.2) -| (a1.east)
              node[pos=0.25, above, font=\scriptsize, red!70!black] {$R^1$};
\draw[reward] (sn.south) -- ++(0,-1.2) -| (aN.east)
              node[pos=0.25, below, font=\scriptsize, red!70!black] {$R^N$};

\end{tikzpicture}
\caption{Марковская игра с $N$ агентами. На каждом шаге каждый агент
  наблюдает общее состояние $s_t$ и выбирает действие. Совместное
  действие $\bm{a}_t$ управляет переходом в $s_{t+1}$, и каждый
  агент получает своё индивидуальное вознаграждение $R^i$.}
\label{fig:markov-game}
\end{figure}

\subsection{Равновесие Нэша}

Поскольку доход каждого агента зависит от стратегий \emph{всех}
агентов, стандартная концепция оптимальности из одноагентного RL
больше не применима. Повышение дохода агента $i$ изменением $\pi^i$ может
потребовать сохранения стратегий других агентов фиксированными\,---\,или может
дать обратный эффект, если они адаптируются. Концепция решения, заимствованная из теории игр,~---
\textbf{равновесие Нэша}\index{Nash equilibrium|textbf}.

\begin{definition}[Равновесие Нэша]
\label{def:nash}
Совместная стратегия $\boldsymbol{\pi}^* = (\pi^{1*}, \ldots, \pi^{N*})$
является \textbf{равновесием Нэша}, если для каждого агента $i$ и каждой
альтернативной стратегии $\tilde\pi^i$:
\[
  J^i(\pi^{i*}, \boldsymbol{\pi}^{-i*}) \;\geq\;
  J^i(\tilde\pi^i, \boldsymbol{\pi}^{-i*}),
\]
где $\boldsymbol{\pi}^{-i*}$ обозначает стратегии всех агентов
кроме $i$, зафиксированные на равновесных значениях.
\end{definition}

Иными словами, равновесие Нэша~--- это совместная стратегия, от которой ни один
агент не имеет стимула отклоняться в одностороннем порядке. Теорема Нэша
гарантирует, что каждая конечная игра имеет хотя бы одно равновесие Нэша
в смешанных стратегиях, но равновесия не обязательно единственны, и нахождение
одного вычислительно сложно в общем случае ($\mathsf{PPAD}$-полная задача для
двухигроковых игр с общей суммой).

\begin{example}[Дилемма заключённого как одношаговая марковская игра]
\label{ex:prisoners}
Два агента одновременно выбирают \textsc{Сотрудничество} ($C$) или
\textsc{Предательство} ($D$). Выигрыши (строка = Агент 1, столбец = Агент 2):
\begin{center}
\small
\begin{tabular}{@{}lcc@{}}
\toprule
 & $C$ & $D$ \\
\midrule
$C$ & $(3,3)$ & $(0,5)$ \\
$D$ & $(5,0)$ & $(1,1)$ \\
\bottomrule
\end{tabular}
\end{center}
Агент 1 всегда предпочитает $D$ независимо от действия Агента 2: если Агент 2
играет $C$, предательство даёт $5 > 3$; если Агент 2 играет $D$, предательство
даёт $1 > 0$. Та же логика применима к Агенту 2. Следовательно, $(D, D)$~---
единственное равновесие Нэша, хотя $(C, C)$ дало бы обоим
агентам более высокий выигрыш $3$. Это противоречие между индивидуальной
рациональностью и коллективным благосостоянием является характерным признаком социальных
дилемм\index{social dilemma} и вновь появится в
Разделе~\ref{sec:communication}.
\end{example}

\subsection{Частичная наблюдаемость и Dec-POMDP}

На практике агенты редко наблюдают полное глобальное состояние. Каждый агент
$i$ вместо этого получает локальное наблюдение $o^i_t$, извлечённое из
функции наблюдения $O^i : \cS \to \Delta(\mathcal{O}^i)$. Получающаяся
постановка~--- \textbf{децентрализованный частично наблюдаемый
MDP}\index{Dec-POMDP} (Dec-POMDP),~--- в худшем случае $\mathsf{NEXP}$-трудна.
Большая часть практических исследований MARL сосредоточена на
ограничении обстановки\,---\,кооперативные команды, общие наблюдения
во время обучения или ограниченная коммуникация,\,---\,чтобы сделать задачу
разрешимой.

% ===========================================================
\section{Кооперативный MARL}
\label{sec:cooperative}
% ===========================================================

\subsection{Кооперативная обстановка}

В полностью кооперативном MARL все агенты разделяют одно командное вознаграждение
$r_t = R(s_t, \bm{a}_t, s_{t+1})$. Совместная цель:
\begin{equation}
\label{eq:cooperative-return}
J(\boldsymbol{\pi}) = \E_{\boldsymbol{\pi}}\!\left[
  \sum_{t=0}^{\infty} \gamma^t r_t \right].
\end{equation}
Несмотря на общую цель, каждый агент должен действовать на основе своего
локального наблюдения и во время исполнения не может общаться или
наблюдать глобальное состояние. Это стандартная обстановка в
многороботной координации, многопользовательских кооперативных играх (например,
микроуправление в StarCraft~II) и многоагентных диалоговых системах.

\subsection{Централизованное обучение с децентрализованным исполнением}

Наиболее влиятельной парадигмой для кооперативного MARL является
\textbf{централизованное обучение с децентрализованным
исполнением}\index{centralised training with decentralised execution|textbf}
(CTDE)\index{CTDE|see{centralised training with decentralised execution}}.
Во время \emph{обучения} центральный критик или микшер имеет доступ к
глобальному состоянию $s_t$ и действиям всех агентов. Во время
\emph{исполнения} каждый агент опирается только на свою локальную историю наблюдений
$\tau^i = (o^i_0, a^i_0, \ldots, o^i_t)$. Привилегия обучения позволяет критику
давать информативный градиентный сигнал; ограничение исполнения гарантирует,
что стратегия применима в обстановках, где центральный
контроллер недоступен.


\begin{figure}[ht]
\centering
\begin{tikzpicture}[
  box/.style={rectangle, draw, rounded corners, minimum width=2.0cm,
              minimum height=0.6cm, font=\small, fill=gray!10, align=center},
  mixer/.style={rectangle, draw, rounded corners, minimum width=2.4cm,
               minimum height=0.7cm, font=\small, fill=blue!12, align=center},
  envbox/.style={rectangle, draw, rounded corners, minimum width=2.4cm,
               minimum height=0.7cm, font=\small, fill=green!8, align=center},
  lossnode/.style={font=\small},
  phase/.style={rectangle, draw, dashed, rounded corners,
                inner sep=10pt},
  arr/.style={->, >=Stealth, thick},
  darr/.style={->, >=Stealth, dashed, gray}
]

%% ============ ФАЗА ОБУЧЕНИЯ ============

\node[box] (q1) at (0,0) {$Q^1(\tau^1, a^1)$};
\node[box] (q2) at (0,-1.2) {$Q^2(\tau^2, a^2)$};
\node[font=\small] (qdots) at (0,-2.1) {$\vdots$};
\node[box] (qN) at (0,-2.8) {$Q^N(\tau^N, a^N)$};

\node[mixer] (mix) at (5.0,-1.2) {\textbf{Микшер} $Q_\mathrm{tot}(s,\bm{a})$};
\node[font=\small] (glob) at (5.0,0.2) {глобальное состояние $s$};
\node[lossnode] (loss) at (7.8,-1.2) {$\mathcal{L}_\mathrm{TD}$};

% Q^i -> Микшер
\draw[arr] (q1.east) -- (mix.west |- q1.east) -- (mix.west);
\draw[arr] (q2.east) -- (mix.west);
\draw[arr] (qN.east) -- (mix.west |- qN.east) -- (mix.west);

% Глобальное состояние -> Микшер
\draw[arr, dashed] (glob.south) -- (mix.north);

% Микшер -> Потери
\draw[arr] (mix.east) -- (loss.west);

\node[phase, fit=(q1)(qN)(glob)(loss),
      label={[font=\small\bfseries]above:Обучение (централизованное)}] (trainbox) {};

%% ============ ФАЗА ИСПОЛНЕНИЯ ============

\node[box] (e1) at (0,-5.2) {Стратегия $\pi^1(\tau^1)$};
\node[box] (e2) at (0,-6.4) {Стратегия $\pi^2(\tau^2)$};
\node[font=\small] (edots) at (0,-7.3) {$\vdots$};
\node[box] (eN) at (0,-8.0) {Стратегия $\pi^N(\tau^N)$};

\node[envbox] (env) at (5.0,-6.4) {Среда};

% Стратегии -> Среда
\draw[arr] (e1.east) -- (env.west |- e1.east) -- (env.west);
\draw[arr] (e2.east) -- (env.west);
\draw[arr] (eN.east) -- (env.west |- eN.east) -- (env.west);

\node[phase, fit=(e1)(eN)(env),
      label={[font=\small\bfseries]above:Исполнение (децентрализованное)}] (execbox) {};

%% ============ ПУНКТИРНЫЕ СВЯЗИ: Q-сети -> Стратегии ============
\draw[darr] (q1.south) -- (e1.north);
\draw[darr] (q2.south) -- (e2.north);
\draw[darr] (qN.south) -- (eN.north);

%% ============ ВЕРТИКАЛЬНЫЙ ТОЧЕЧНЫЙ РАЗДЕЛИТЕЛЬ ============
\coordinate (septop) at (2.5, 0.9);
\coordinate (sepbot) at (2.5,-8.7);
\draw[dotted, thick, gray!60] (septop) -- (sepbot);

\end{tikzpicture}
\caption{Архитектура CTDE. Во время обучения \emph{микшер}
  (например, QMIX) объединяет сети полезности $Q^i$ каждого агента в
  совместный $Q_\mathrm{tot}$ с использованием глобального состояния, и градиенты
  распространяются через все $Q^i$. Во время исполнения централизованный микшер
  отбрасывается; каждый агент действует жадно относительно своего $Q^i$,
  используя только свою локальную историю наблюдений $\tau^i$.}
\label{fig:ctde}
\end{figure}

\subsection{QMIX}

\textbf{QMIX}\index{QMIX|textbf} (Rashid et al., 2018)~--- канонический алгоритм CTDE
для кооперативного MARL на основе функции ценности. Каждый
агент $i$ поддерживает сеть полезности
$Q^i_\theta(\tau^i, a^i)$, обусловленную только на локальной
истории наблюдений. Центральная \emph{сеть-микшер} объединяет эти
полезности в совместную функцию ценности действий:
\begin{equation}
\label{eq:qmix}
Q_\mathrm{tot}(s, \bm{a}) = f_\phi\!\bigl(Q^1(\tau^1, a^1), \ldots,
  Q^N(\tau^N, a^N),\; s\bigr).
\end{equation}
Микшер использует гиперсети\index{hypernetwork}, обусловленные на
глобальном состоянии $s$, для получения неотрицательных весов, что гарантирует
\textbf{ограничение монотонности}\index{QMIX!monotonicity}:
\[
  \frac{\partial Q_\mathrm{tot}}{\partial Q^i} \geq 0
  \quad \text{для всех } i.
\]
Это ограничение является ключевым инсайтом. Оно гарантирует
свойство \textbf{Индивидуально-Глобального Максимума}\index{IGM condition} (IGM):
\[
  \argmax_{\bm{a}}\, Q_\mathrm{tot}(s, \bm{a})
  \;=\; \bigl(\argmax_{a^1} Q^1(\tau^1, a^1),\; \ldots,\;
               \argmax_{a^N} Q^N(\tau^N, a^N)\bigr).
\]
Максимизация $Q_\mathrm{tot}$ эквивалентна тому, что каждый агент
жадно выбирает своё лучшее действие\,---\,поиск по совместным действиям
во время исполнения не нужен. Вся система обучается end-to-end
со стандартной функцией потерь TD:
\begin{equation}
\label{eq:qmix-loss}
\mathcal{L}(\theta, \phi) = \E\!\left[
  \bigl(r + \gamma\, \max_{\bm{a}'} Q_\mathrm{tot}(s', \bm{a}')
  - Q_\mathrm{tot}(s, \bm{a})\bigr)^2
\right].
\end{equation}
QMIX достигает лучших результатов в области на тесте StarCraft
Multi-Agent Challenge (SMAC)\index{SMAC}~--- наборе задач кооперативного
микроуправления в StarCraft~II.

\subsection{MAPPO}

Тогда как QMIX работает в парадигме функции ценности, многие кооперативные
задачи выигрывают от методов градиента стратегии, поскольку они естественно обрабатывают
непрерывные действия и проще в настройке. \textbf{MAPPO}
\index{MAPPO|textbf} (Yu et al., 2022) применяет Proximal Policy
Optimisation (PPO; Глава~\ref{ch:actor-critic}) в фреймворке CTDE.
Каждый агент $i$ поддерживает отдельного актора $\pi^i_\theta(\cdot \mid
\tau^i)$, обусловленного на его локальной истории наблюдений. \emph{Общий
централизованный критик} $V_\phi(s)$ обусловлен на глобальном состоянии и
производит baseline для снижения дисперсии:
\begin{equation}
\label{eq:mappo-objective}
\mathcal{L}_\mathrm{MAPPO}(\theta^i) =
\E\!\left[\min\!\bigl(r^i_t(\theta^i)\,\hat{A}_t,\;
  \mathrm{clip}(r^i_t(\theta^i), 1-\epsilon, 1+\epsilon)\,\hat{A}_t
\bigr)\right],
\end{equation}
где $r^i_t(\theta^i) = \pi^i_\theta(a^i_t \mid \tau^i_t) /
\pi^i_{\mathrm{old}}(a^i_t \mid \tau^i_t)$~--- коэффициент вероятностей,
а $\hat{A}_t$~--- преимущество, оцениваемое централизованным критиком.
Ключевой практической находкой является то, что \emph{совместное использование параметров}\index{parameter
sharing}\,---\,инициализация всех агентов из одной сети актора
при обусловлении на признаке индекса агента\,---\,резко
улучшает эффективность по образцам, позволяя агентам разделять градиенты.

\begin{remark}[Совместное использование параметров против гетерогенных агентов]
Совместное использование параметров мощно, когда агенты играют симметричные роли, но
может быть неуместным для команд из разнородных специалистов (например,
целителя и танка в игре). В этой обстановке отдельные сети с
общим экстрактором признаков часто работают лучше всего.
\end{remark}

% ===========================================================
\section{Конкурентные и смешанные обстановки}
\label{sec:competitive}
% ===========================================================

\subsection{Игры с нулевой суммой и самоигра}

В двухигроковой игре с нулевой суммой $R^1 = -R^2$, так что выигрыш одного агента~---
ровно проигрыш другого. Равновесие Нэша игры с нулевой суммой по Маркову
соответствует стратегии минимакса: агент 1 максимизирует, а агент
2 минимизирует одно и то же накопленное вознаграждение. Формально минимаксное значение
$v^*$ удовлетворяет:
\[
  v^* = \max_{\pi^1} \min_{\pi^2}\; J^1(\pi^1, \pi^2)
       = \min_{\pi^2} \max_{\pi^1}\; J^1(\pi^1, \pi^2),
\]
где равенство (теорема о минимаксе фон Неймана)\index{minimax theorem}
выполняется, когда стратегии могут быть смешанными (стохастическими).

Обучение агента для оптимальной игры в игре с нулевой суммой сталкивается с
непосредственной проблемой начальной загрузки: чтобы научиться побеждать сильных противников,
агенту нужны сильные противники для обучения. \textbf{Самоигра}
\index{self-play|textbf} решает это, заставляя агента играть против
прошлых версий самого себя. Канонический цикл обучения:
\begin{enumerate}[leftmargin=*]
  \item Поддерживать популяцию (или один снимок) прошлых контрольных точек стратегии.
  \item Выбрать противника из популяции.
  \item Собрать развёртки; обновить текущую стратегию против этого
    противника.
  \item Периодически добавлять текущую стратегию в популяцию.
\end{enumerate}
Самоигра гарантирует, что распределение противников всегда имеет
уровень сложности, соразмерный текущей стратегии, обеспечивая
естественный учебный план.

\subsection{AlphaGo и AlphaZero}

Наиболее известное применение самоигры~--- программы AlphaGo и
AlphaZero\index{AlphaZero|textbf} (Silver et al., 2016;
2017). Го\index{Go (game)} долго сопротивлялось решениям ИИ, поскольку его
игровое дерево имеет коэффициент ветвления около $250$ и глубину более
$200$, делая исчерпывающий поиск невозможным.

AlphaGo объединял три компонента: (1) сеть ценности $v_\theta(s)$,
обученную регрессией на результатах самоигры; (2) сеть стратегии
$\pi_\theta(a \mid s)$, обученную обучением с учителем на партиях людей
и затем уточнённую градиентом стратегии RL; и (3) поиск по дереву Монте-Карло
(MCTS)\index{MCTS}, использующий обе сети для управления развёртками.
На каждом узле MCTS шаг выбора балансирует исследование и
использование через критерий PUCT\index{PUCT}:
\begin{equation}
\label{eq:puct}
a^* = \argmax_{a}\!\left[
  Q(s, a) + c_\mathrm{puct}\, \pi_\theta(a \mid s)
  \frac{\sqrt{N(s)}}{1 + N(s, a)}
\right],
\end{equation}
где $Q(s,a)$~--- эмпирическое среднее значение из прошлых развёрток,
$N(s)$~--- счётчик посещений узла $s$, а $N(s,a)$~--- счётчик
ребра $(s,a)$.

AlphaZero\index{AlphaZero} убрал начальную загрузку через обучение с учителем:
самоигра вместе с MCTS, даже при старте с \emph{random play},
оказались достаточны для достижения сверхчеловеческой производительности в Го, шахматах
и сёги за несколько дней обучения. Успех AlphaZero
показал, что при наличии идеального симулятора самоигра может открыть
экспертные стратегии полностью с нуля\,---\,знаковый результат как для
ИИ, так и для MARL.

\subsection{Нестабильность конкурентного обучения и нетранзитивные стратегии}

Наивная самоигра не лишена патологий. В играх с
\textbf{нетранзитивными стратегическими отношениями}\index{non-transitive strategies}\,---\,где стратегия
$A$ бьёт $B$, $B$ бьёт $C$, но $C$ бьёт $A$ (динамика «камень-ножницы-бумага»)\,---\,самоигра может циклически переходить вместо сходимости
к равновесию Нэша. Стратегия, обученная против версии $k$ самой себя,
может быть сильной против $k$, но хрупкой против версии $k-5$~--- проблема,
известная как \textbf{забывание}\index{forgetting!in self-play} или
стратегический коллапс.

Широко используются два решения. \textbf{Фиктивная
самоигра}\index{fictitious self-play} обучает агента против
\emph{среднего} всех прошлых стратегий противника, а не только самого последнего
снимка. В двухигроковых играх с нулевой суммой с конечными пространствами действий
фиктивная самоигра сходится к равновесию Нэша.
\textbf{Популяционное обучение}\index{population-based training}
(PBT; Jaderberg et al., 2017) поддерживает разнообразную лигу агентов с
различными стратегиями; каждый агент в популяции обучается против
смеси других, предотвращая устранение разнообразия любой одной доминирующей стратегией.

% ===========================================================
\section{Коммуникация и возникающее поведение}
\label{sec:communication}
% ===========================================================

\subsection{Обученные протоколы коммуникации}

Когда агенты могут обмениваться сообщениями до или во время исполнения задачи,
эффективное пространство стратегий команды резко расширяется. Но
что им говорить? Жёсткое кодирование протокола коммуникации требует
экспертизы в предметной области и плохо масштабируется. Более принципиальный подход~---
рассматривать сообщения как \emph{непрерывные действия} и позволять протоколам коммуникации
возникать через end-to-end обучение.

\textbf{CommNet}\index{CommNet|textbf} (Sukhbaatar et al., 2016)~---
одна из первых архитектур для дифференцируемой многоагентной
коммуникации. На каждом шаге каждый агент $i$ транслирует
непрерывный вектор $m^i_t \in \R^d$; агрегированное сообщение
$\bar{m}^i_t = \frac{1}{N-1}\sum_{j \neq i} m^j_t$ конкатенируется
с наблюдением агента и подаётся в LSTM\index{LSTM}, чей
выход управляет как следующим действием, так и следующим сообщением:
\[
  (a^i_t, m^i_{t+1}) = f_\theta\!\bigl(h^i_t, o^i_t, \bar{m}^i_t\bigr).
\]
Поскольку канал коммуникации дифференцируем, вся
система обучается end-to-end со стандартным обратным распространением через
время. CommNet показал, что агенты могут спонтанно обучиться
координации на кооперативных задачах (например, управление трафиком на перекрёстках)
без какого-либо надзора над содержанием коммуникации.

\textbf{TarMAC}\index{TarMAC|textbf} (Das et al., 2019) расширил
CommNet, введя \emph{целевую коммуникацию}: каждый агент
производит ключ $k^i$ и значение $v^i$, а агрегированное сообщение,
полученное агентом $j$,~--- взвешенная по вниманию сумма:
\[
  m^j = \sum_{i \neq j} \sigma(k^i \cdot q^j)\, v^i,
\]
где $q^j$~--- вектор запроса агента $j$. Механизм внимания
позволяет агентам избирательно читать наиболее релевантные сообщения для их
текущей ситуации, улучшая как производительность, так и
интерпретируемость обученной коммуникации.

\subsection{Возникающая коммуникация}

Поразительный результат литературы по возникающей коммуникации состоит в том, что
агенты могут развивать \emph{композиционные языки}\index{emergent
language!compositionality}: короткие дискретные символы, комбинации которых
кодируют структурированные смыслы. В референтных играх\index{referential game}
(например, один агент видит объект и должен его описать; другой агент
должен идентифицировать его среди отвлекающих объектов) обученные агенты часто развивают
словари, в которых разные символы соответствуют разным атрибутам объекта,
а новые комбинации символов обозначают новые объекты.
Эти находки параллельны изучению эволюции языка и поднимают
вопрос о том, могут ли искусственные возникающие языки стать основой
для полезных коммуникативных навыков для агентов реального мира.

\subsection{Социальные дилеммы в многоагентном обучении}

Дилемма заключённого из Примера~\ref{ex:prisoners}~--- одношаговая
игра. В \emph{повторяющихся} играх кооперация может возникать спонтанно,
поскольку агенты имеют возможность \emph{наказывать} предателей в будущих
раундах. Знаменитые турниры Аксельрода показали, что простая
стратегия \textsc{Тит за тат}\index{Tit-for-Tat} (кооперировать
изначально; на каждом последующем ходу копировать последнее действие противника)
превосходит сложные стратегии в повторяющихся дилеммах заключённого.

В MARL социальные дилеммы\index{social dilemma} формализуются как
среды, где равновесие Нэша является Парето-субоптимальным: рациональное
индивидуальное поведение ведёт к коллективно плохим результатам.
Фреймворк последовательной социальной дилеммы (SSD)\index{sequential social dilemma}
(Leibo et al., 2017) встраивает такие дилеммы в пространственные
среды. Например, в игре \emph{Gathering} два агента
собирают яблоки и могут стрелять «штрафными лучами» для временного отключения
друг друга. Взаимная стрельба отнимает время, но каждый агент индивидуально
стимулируется выстрелить первым. Агенты RL, обученные независимым
обучением, сходятся к взаимному предательству; агенты, обученные с просоциальным
формированием вознаграждения или явной коммуникацией, учатся кооперировать.

% ===========================================================
\section{Мультиагентный RL для LLM}
\label{sec:marl-llms}
% ===========================================================

В предыдущих разделах агенты рассматривались как нейронные сети, отображающие
наблюдения в действия в среде с ясными состояниями и сигналами вознаграждения.
Большие языковые модели (LLM)\index{large language model}
вводят качественно новый вид агента: один, работающий в
пространстве естественного языка, способный рассуждать через промежуточные шаги и
делегировать подзадачи другим экземплярам LLM. Концепции MARL применяются
естественно, и растущий массив работ использует их для улучшения
качества рассуждений, генерации кода и выполнения задач.

\subsection{Дебаты}

\textbf{Дебаты ИИ}\index{debate (MARL)} (Irving et al., 2018)~---
многоагентный протокол для получения надёжных выходов от
сильной системы ИИ. Два языковых агента отстаивают противоположные стороны
утверждения; третий агент (или судья-человек) оценивает
обмен и объявляет победителя. Ключевое теоретическое утверждение состоит в том, что
для достаточно способного судьи оптимальная стратегия каждого дебатёра~---
\emph{говорить правду}: ложь всегда может быть разоблачена
противником, поэтому лучшая стратегия в равновесии дебатов~--- честная аргументация.

На практике протоколы в стиле дебатов реализуются как пайплайн из двух LLM:
\begin{enumerate}[leftmargin=*]
  \item LLM-\textbf{предлагающий} генерирует утверждение или ответ.
  \item LLM-\textbf{критик}, получив тот же промпт и выход
    предлагающего, генерирует контраргумент.
  \item LLM-\textbf{судья} (или та же модель с промптом судьи)
    оценивает обе позиции и выбирает лучший ответ.
\end{enumerate}
Оценка судьи может использоваться как сигнал вознаграждения для дообучения как
предлагающего, так и критика через RL (Раздел~\ref{sec:rlvr-ch15} в
Главе~\ref{ch:reasoning}). Это обеспечивает \emph{масштабируемый надзор}:
система может улучшаться за пределами прямых возможностей судьи,
поскольку судье проще оценить аргумент, чем
генерировать правильный ответ с нуля.

\subsection{Многоагентная генерация кода}

Генерация кода~--- естественная обстановка для многоагентных LLM, поскольку
код допускает автоматические верифицируемые вознаграждения (наборы тестов), а
задача разбивается на подзадачи (спецификация, реализация,
тестирование, документация), которые легко назначать специализированным агентам.

Представительный пайплайн (Chen et al., 2023; Du et al., 2023) работает
следующим образом:
\begin{enumerate}[leftmargin=*]
  \item Агент-\textbf{планировщик} разбивает запрос пользователя на
    последовательность подзадач.
  \item Несколько агентов-\textbf{программистов} независимо генерируют
    кандидатные реализации каждой подзадачи.
  \item Агент-\textbf{верификатор} запускает наборы тестов на каждом кандидате
    и возвращает обратную связь об успехе/неудаче.
  \item Агент-\textbf{ранкер}, обученный с предпочтительным RL, выбирает
    реализацию с наивысшим рангом и при необходимости исправляет неудачи.
\end{enumerate}
Эмпирически разнообразные предложения от нескольких независимых программистов
стабильно превосходят любого одного программиста, даже когда отдельные
программисты являются идентичными моделями: разнообразие в случайных семенах, промптах или
настройках температуры вводит достаточную вариацию, чтобы верификатор
мог выявлять действительно лучшие решения.

\subsection{LLM-агенты, сотрудничающие над задачами}

Помимо кода, мультиагентные LLM-системы применялись к исследовательской
помощи, игровым задачам и программной инженерии. \textbf{ChatDev}
(Qian et al., 2023) структурирует программную команду как набор LLM-агентов,
играющих разные роли (генеральный директор, технический директор, программист, тестировщик, автор документации),
которые общаются через модерируемый чат-канал для завершения программного
проекта от начала до конца. Каждая роль определяется системным промптом и
ролевой историей; протокол модерации обеспечивает очерёдность ходов
и фокус на задаче.

\textbf{AutoGen}\index{AutoGen} (Wu et al., 2023) предоставляет более
гибкий фреймворк: любой LLM может быть обёрнут как \emph{агент} с
именем, системным промптом и при желании интерфейсом использования инструментов. Агенты
общаются, передавая сообщения в общем потоке; агент, представляющий
человека, может вводить обратную связь от пользователя. Фреймворк поддерживает циклы, ветвящиеся
разговоры и вложенные группы агентов, позволяя как кооперативные, так и
состязательные паттерны.

\subsection{Agentic Workflows со специализированными агентами}

Наиболее практически влиятельным паттерном мультиагентного LLM на момент
написания является паттерн \textbf{оркестратор\,---\,специалист}.
Мощный универсальный LLM (\emph{оркестратор}) получает сложную
задачу пользователя, разбивает её на подзадачи и направляет каждую подзадачу
\emph{специализированному} агенту, оптимизированному для конкретного навыка:
\begin{description}[leftmargin=!, labelwidth=3.5cm]
  \item[Агент веб-поиска] извлекает и резюмирует актуальную информацию.
  \item[Агент кода] пишет и выполняет код в изолированном
    интерпретаторе.
  \item[Агент памяти] ведёт постоянное хранилище ключ-значение фактов,
    извлечённых из предыдущих ходов.
  \item[Агент-критик] проверяет план оркестратора и отмечает
    несоответствия перед исполнением.
\end{description}
Вознаграждения для дообучения RL оркестратора могут поступать от
end-to-end решения задачи (например, совпадает ли финальный ответ с
эталонным решением?) или от промежуточных подцелей, проверяемых
специализированными агентами. Этот замкнутый цикл между оркестрацией и
исполнением специалистом отражает парадигму CTDE: оркестратор
действует как централизованный планировщик во время обучения, в то время как каждый специалист
выполняет свою локальную стратегию независимо.

Связь с MARL прямая. Каждый специалист~--- это стратегия
$\pi^i$; оркестратор определяет распределение задач над совместными
последовательностями действий; а возникающее поведение ансамбля~--- это
траектория в пространстве всех возможных многоагентных взаимодействий.
Формальное обучение этой системы с целями MARL\,---\,а не
дообучение каждого компонента в отдельности,\,---\,является активным рубежом исследований.

% ===========================================================
\section{Проблемы мультиагентного RL}
\label{sec:marl-challenges}
% ===========================================================

\subsection{Нестационарность}

Наиболее фундаментальная проблема MARL состоит в том, что среда каждого агента
включает всех других агентов, которые сами
обучаются. С точки зрения агента $i$ закон перехода и
функция вознаграждения \emph{нестационарны}\index{non-stationarity!in MARL|textbf}:
мир, который он испытывает сегодня, отличается от мира, который он будет
испытывать завтра, поскольку его товарищи по команде или противники обновили
свои стратегии. Это нарушает допущение стационарности, лежащее в основе
доказательств сходимости для Q-обучения и других одноагентных
алгоритмов.

Формально агент $i$ наблюдает эффективный закон перехода:
\[
  \tilde{\cT}^i(s' \mid s, a^i)
  = \sum_{\bm{a}^{-i}} \cT(s' \mid s, a^i, \bm{a}^{-i})
    \prod_{j \neq i} \pi^j(a^j \mid s),
\]
который меняется каждый раз, когда меняется любая $\pi^j$ ($j \neq i$). Независимое
Q-обучение\index{independent Q-learning} игнорирует эту нестационарность
и часто не сходится на практике, особенно в конкурентных
обстановках. Алгоритмы, такие как QMIX и MAPPO, частично смягчают
нестационарность путём совместного использования информации о градиентах между агентами, но
гарантии сходимости остаются слабыми в общем случае.

\subsection{Отнесение заслуг в командах}

Даже при общем сигнале вознаграждения командное отнесение заслуг
сложно\index{credit assignment!in MARL}: если команда добивается успеха, чьё
действие было ответственным? Наивные подходы равномерно распределяют всё вознаграждение
между всеми агентами, что даёт шумные градиенты и препятствует
специализации. Эвристика \textbf{разностного вознаграждения}\index{difference reward}
присваивает агенту $i$ кредит:
\[
  d^i_t = r_t - r_t^{-i},
\]
где $r_t^{-i}$~--- вознаграждение, которое \emph{было бы} получено,
если бы агент $i$ выполнил какое-то дефолтное или контрфактическое действие,
при зафиксированных других агентах. Вычисление $r_t^{-i}$ требует либо
обученной модели мира, либо дополнительных развёрток, что делает его дорогостоящим.
Ограничение монотонности QMIX обеспечивает неявный механизм отнесения заслуг:
градиент $Q_\mathrm{tot}$ по $Q^i$
говорит агенту $i$, насколько его функция полезности вносит вклад в команду.

\subsection{Масштабируемость}

Совместное пространство действий $|\cA| = \prod_i |\cA^i|$ растёт экспоненциально
с числом агентов. Для $N = 10$ агентов каждый с $|\cA^i| = 5$
действиями совместное пространство имеет $5^{10} \approx 10^7$ элементов. Даже
представление $Q_\mathrm{tot}(s, \bm{a})$ в виде таблицы неосуществимо;
факторизация CTDE в QMIX решает это, записывая
$Q_\mathrm{tot}$ как функцию от полезностей каждого агента, уменьшая
сложность с экспоненциальной до линейной по $N$.

Методы градиента стратегии лучше масштабируются, поскольку они работают с
маргинальными стратегиями $\pi^i(a^i \mid \tau^i)$, а не с совместной
стратегией. С совместным использованием параметров одна сеть обслуживает всех агентов,
поэтому число параметров не растёт с $N$. Тем не менее
время обучения по-прежнему масштабируется с числом агентов, поскольку нужно больше
данных развёрток, а межагентные зависимости требуют больших
буферов воспроизведения и мини-батчей.

\subsection{Формирование вознаграждения для кооперации}

В играх со смешанными мотивами агенты могут иметь индивидуально рациональные
стимулы, конфликтующие с командной производительностью. Формирование вознаграждения
\index{reward shaping!in MARL}\,---\,дополнение истинного вознаграждения
потенциально-основанным бонусом,\,---\,может подтолкнуть агентов к
кооперативному поведению без нарушения равновесия Нэша.
Однако плохое формирование вознаграждений может привносить ложные равновесия
или делать агентов безразличными к истинной цели.

\textbf{Внутренняя социальная мотивация}\index{intrinsic motivation!social}~---
многообещающая альтернатива: агентам даётся внутреннее вознаграждение
за \emph{влияние} на поведение других (вознаграждение за социальное влияние;
Jaques et al., 2019) или за достижение causal empowerment
над средой. Эти внутренние вознаграждения поощряют обмен информацией
без необходимости в вручную разработанных протоколах коммуникации.

\subsection{Эффективность по образцам}

MARL требует много образцов. Не только каждый агент должен исследовать своё
пространство действий, но и совместная проблема исследования экспоненциально
больше. Координированное исследование\,---\,где агенты активно
избегают избыточного исследования одних и тех же областей состояний,\,---\,уменьшает
избыточность, но требует дополнительных затрат на координацию. Офлайн
MARL\index{offline MARL}, который учится из заранее собранных наборов данных
без онлайн-взаимодействия~--- это появляющийся подход, переносящий
недавние успехи офлайн одноагентного RL
(Глава~\ref{ch:open}) в многоагентную обстановку.

В таблице ниже обобщены основные алгоритмические семейства,
их обстановки и ключевые свойства.

\begin{center}
\footnotesize
\begin{tabular}{@{}p{1.5cm}p{2cm}p{2.1cm}p{2.1cm}p{3.9cm}@{}}
\toprule
\textbf{Алгоритм} & \textbf{Обстановка} & \textbf{Парадигма} & \textbf{Пространство действий} & \textbf{Ключевая идея} \\
\midrule
IQL & Коо./Кон. & Децентрализ. & Дискретное & Независимое Q-обучение \\
QMIX & Кооперат. & CTDE & Дискретное & Монотонный микшер \\
QPLEX & Кооперат. & CTDE & Дискретное & Микшер с двойным вниманием \\
MAPPO & Кооперат. & CTDE & Непр./Диск. & Общий центральный критик \\
MADDPG & Смешанный & CTDE & Непрерывное & Центральный критик каждого агента \\
AlphaZero & Конкурент. & Самоигра & Дискретное & MCTS + сети стратегии/ценности \\
PSRO & Смешанный & Популяция & Дискретное & Nash resp.\ oracle + mixture \\
CommNet & Кооперат. & Канал связи & Непрерывное & Сообщения среднего поля \\
TarMAC & Кооперат. & Канал связи & Непрерывное & Целевые сообщения с вниманием \\
\bottomrule
\end{tabular}
\end{center}

% ===========================================================
\section{Итоги}
\label{sec:marl-summary}
% ===========================================================

Мультиагентное обучение с подкреплением расширяет одноагентный формализм MDP
на $N$ взаимодействующих агентов, каждый со своей стратегией,
наблюдением и вознаграждением. Центральная концепция решения~---
равновесие Нэша: совместная стратегия, от которой ни один агент не может выгодно
отклониться в одностороннем порядке. Нахождение равновесий Нэша сложно в общем случае,
но структура игры\,---\,кооперация, конкуренция с нулевой суммой,
факторизованная динамика,\,---\,обеспечивает разрешимые алгоритмы.

Для кооперативных задач парадигма CTDE достигает лучшего из
обоих миров: централизованная информация о градиентах во время обучения и
полностью децентрализованное исполнение. QMIX использует монотонную
факторизацию для обеспечения глобальной оптимальности жадного выбора действий каждым агентом.
MAPPO применяет PPO с централизованным value baseline и
совместным использованием параметров, масштабируясь на большие команды с непрерывными действиями.

В конкурентных обстановках самоигра обеспечивает автоматический учебный план.
AlphaZero показывает, что даже при старте с random play
сочетания самоигры и MCTS достаточно для сверхчеловеческой производительности в
играх с полной информацией. Нетранзитивность и стратегическое забывание
мотивируют популяционные методы, поддерживающие разнообразие стратегий.

Коммуникация открывает дополнительный канал для кооперации.
CommNet и TarMAC демонстрируют, что end-to-end обучение может
производить эффективные протоколы обмена сообщениями без какого-либо надзора
над содержанием сообщений. В социальных дилеммах, где равновесие Нэша
является Парето-субоптимальным, повторяющееся взаимодействие и просоциальные
вознаграждения могут поддерживать кооперативное поведение.

LLM~--- последняя и, пожалуй, наиболее значимая арена для идей MARL.
Дебаты, многоагентная генерация кода, пайплайны оркестратор-специалист
и системы вроде AutoGen и ChatDev реализуют
многоагентную динамику в языке. Применение формальных целей обучения MARL
к этим системам\,---\,а не дообучение компонентов
по отдельности,\,---\,является одним из наиболее захватывающих открытых направлений на
рубеже исследований RL.

Ключевые проблемы\,---\,нестационарность, отнесение заслуг,
масштабируемость, формирование вознаграждения и эффективность по образцам,\,---\,взаимосвязаны, и ни одна полностью не решена. Прогресс в каждой, вероятно,
потребует новой теории и новой инженерии, что делает MARL одной из
наиболее активных и значимых подобластей машинного обучения.

% ===========================================================
\section*{Упражнения}
\addcontentsline{toc}{section}{Упражнения}
% ===========================================================

\begin{exercise}
\label{ex:markov-game-reduction}
Покажите, что марковская игра с $N=1$ сводится к стандартному MDP.
Затем покажите, что если все агенты разделяют одно вознаграждение $R^1 = \cdots =
R^N$ и все агенты действуют одинаково (т.е. $\pi^1 = \cdots = \pi^N$),
совместный доход из Уравнения~\eqref{eq:cooperative-return} равен
одноагентному доходу из Главы~\ref{ch:mdp}. Какие дополнительные
условия нужны для того, чтобы градиент стратегии с общим вознаграждением равнялся
одноагентному градиенту стратегии?
\end{exercise}

\begin{exercise}
\label{ex:nash-prisoners}
Рассмотрим дилемму заключённого из Примера~\ref{ex:prisoners}.
\begin{enumerate}[label=(\alph*)]
  \item Убедитесь, что $(D, D)$~--- равновесие Нэша, проверив
    определение: ни один игрок не может улучшить свой выигрыш, перейдя к $C$.
  \item Покажите, что $(C, C)$~--- \emph{не} равновесие Нэша.
  \item Теперь предположим, что игра проводится $T = 100$ раундов и
    агенты максимизируют суммарное недисконтированное вознаграждение. Является ли
    равновесием Нэша для обоих агентов всегда кооперировать? Обоснуйте.
  \item Что меняется, если горизонт бесконечен с коэффициентом дисконтирования
    $\gamma = 0.9$?
\end{enumerate}
\end{exercise}

\begin{exercise}
\label{ex:igm-proof}
Свойство IGM утверждает, что
\[
\argmax_{\bm{a}} Q_\mathrm{tot}(s, \bm{a}) = \bigl(\argmax_{a^1}
Q^1(\tau^1, a^1),\, \ldots,\, \argmax_{a^N} Q^N(\tau^N, a^N)\bigr).
\]
\begin{enumerate}[label=(\alph*)]
  \item Предположим, что $Q_\mathrm{tot}(s, \bm{a}) = \sum_i Q^i(\tau^i,
    a^i)$. Докажите, что свойство IGM выполняется.
  \item Выполняется ли свойство IGM, если
    $Q_\mathrm{tot}(s, \bm{a}) = \prod_i Q^i(\tau^i, a^i)$ и
    все $Q^i > 0$? Докажите или приведите контрпример.
  \item Объясните в одном абзаце, почему неотрицательные смешивающие
    веса QMIX (из его гиперсети) достаточны, но не необходимы
    для IGM.
\end{enumerate}
\end{exercise}

\begin{exercise}
\label{ex:self-play-rps}
Рассмотрим игру с тремя стратегиями: чистыми стратегиями Камень (R), Ножницы
(P), Бумага (S) и матрицей выигрышей:
\[
M = \begin{pmatrix} 0 & -1 & 1 \\ 1 & 0 & -1 \\ -1 & 1 & 0
\end{pmatrix}
\]
где $M_{ij}$~--- выигрыш игрока по строке при игре стратегией $i$
против игрока по столбцу со стратегией $j$.
\begin{enumerate}[label=(\alph*)]
  \item Найдите единственное равновесие Нэша этой игры.
  \item Симулируйте $1\,000$ шагов градиентного подъёма по $\pi^1$, пока
    $\pi^2$ фиксирована в равновесии Нэша. Сходится ли $\pi^1$?
  \item Теперь симулируйте $1\,000$ шагов одновременного градиентного подъёма
    для обоих игроков. Опишите траекторию в пространстве смешанных
    стратегий. Сходится ли она к равновесию Нэша?
  \item Интуитивно объясните, почему эта игра вызывает циклирование наивной самоигры.
\end{enumerate}
\end{exercise}

\begin{exercise}
\label{ex:credit-assignment}
Кооперативная команда из $N = 4$ роботов получает общее вознаграждение $r =
10$, когда все четверо достигают цели одновременно, и $r = 0$ иначе.
После успешного эпизода независимое Q-обучение равномерно распределяет
кредит $10/4 = 2.5$ каждому агенту.
\begin{enumerate}[label=(\alph*)]
  \item Предположим, что Робот 3 имел вероятность $0.5$ достичь цели
    независимо от своего действия. Опишите, как равномерное распределение кредита
    искажает сигнал обучения для Роботов 1, 2 и 4.
  \item Определите разностное вознаграждение для Робота $i$ с помощью контрфактического
    сценария, где Робот $i$ стоит на месте, пока другие следуют наблюдаемой стратегии.
    Вычислите это вознаграждение для Робота 3.
  \item В чём преимущество разностного вознаграждения перед равномерным
    кредитом? В чём его главный вычислительный недостаток?
\end{enumerate}
\end{exercise}

\begin{exercise}
\label{ex:ctde-communication}
Разработайте алгоритм CTDE для команды из $N$ дронов, обследующих
зону поиска и спасения. Каждый дрон наблюдает локальную сетку $5 \times 5$; глобальная
карта $50 \times 50$. Агенты не могут общаться во время тестирования.
\begin{enumerate}[label=(\alph*)]
  \item Определите пространство наблюдений $\mathcal{O}^i$, пространство действий
    $\mathcal{A}^i$ и подходящую командную функцию вознаграждения $R$.
  \item Предложите архитектуру централизованного критика, использующую
    полную карту $50 \times 50$ во время обучения. Какие признаки вы бы
    конкатенировали наряду с глобальным состоянием?
  \item Обоснуйте, уместно ли здесь совместное использование параметров, и
    как бы вы обусловили актора каждого агента на его уникальном
    идентификаторе.
  \item Определите два способа возникновения нестабильности обучения и
    предложите меры по их устранению.
\end{enumerate}
\end{exercise}

\begin{exercise}
\label{ex:debate-reward}
Реализуйте упрощённый протокол дебатов для задачи closed-book
question answering.
\begin{enumerate}[label=(\alph*)]
  \item Определите состояние, действие и вознаграждение для каждого из
    агентов-предлагающего, критика и судьи в однораундовых дебатах.
  \item Пусть судья присваивает бинарное вознаграждение $r \in \{0, 1\}$
    (верно / неверно). Запишите обновление градиента стратегии REINFORCE
    для предлагающего, предполагая, что стратегия критика
    зафиксирована. Высока ли дисперсия этого оценщика?
    Почему?
  \item Предложите baseline для предлагающего, снижающий дисперсию
    без введения смещения. Конкретно опишите, как вы бы её оценили.
  \item При каких условиях равновесие дебатов гарантирует
    правдивое поведение предлагающего? Сформулируйте необходимые
    допущения о способностях судьи.
\end{enumerate}
\end{exercise}

\begin{exercise}
\label{ex:emergent-language}
В референтной игре агент-говорящий наблюдает изображение, извлечённое из
множества $K$ объектов, и должен произвести дискретное сообщение
$m \in \{1, \ldots, V\}^L$ (размер словаря $V$, длина сообщения $L$).
Слушатель наблюдает $m$ и должен идентифицировать правильный объект из $K$
отвлекающих объектов.
\begin{enumerate}[label=(\alph*)]
  \item Какое максимальное число различных объектов может быть
    закодировано языком со словарём $V$ и длиной $L$?
  \item Предположим, что оба агента обучаются end-to-end с REINFORCE,
    и вознаграждение равно $+1$ за верную идентификацию, $0$ иначе.
    Запишите совместное обновление градиента стратегии.
  \item Опишите одну метрику для измерения композиционности
    возникшего языка (например, топографическое сходство). Определите её
    формально.
  \item Почему возникший язык может не быть композиционным,
    даже когда композиционность была бы полезна?
\end{enumerate}
\end{exercise}
